\begin{proposition}[A noncontracted horizontal cusp has zero Euler number]\hypertarget{tag-01BB}{}\label{cdp:st:A.5.20}
\StatementRecord{cdp:st:A.5.20}{proved}{01BB}{\SeriesRef{cdp:st:A.5.14}, \SeriesRef{cdp:st:A.5.5}, \SeriesRef{cdp:st:C2d}, \SeriesRef{cdp:st:A.5.18}}

Let $f:Y\to C_0$ be a proper holomorphic map from a smooth
compact complex threefold with $b_2(Y)=0$, with connected
fibres and smooth general fibre a complex two-torus.
Let $mS$ be a full fibre with irreducible reduced support,
$\eta:T\to S$ its normalization and $\pi:Z\to T$ an
irrational ruled resolution with ruling $p:Z\to B$.
Suppose every exceptional curve of $\pi$ is vertical
and the resolution conductor is $2C+V$, where $C$
is a noncontracted section and $V$ is vertical.
If the generic transversal at $D=\eta\pi(C)$ is an
ordinary cusp, then $D$ and $C$ are smooth elliptic
curves and $C\to D$ is an isomorphism. The cusp type
is constant along $D$. The normalization $T$ is a
smooth minimal ruled surface over this elliptic curve,
its conductor is twice a section, and $e(S)=0$.
\end{proposition}


\begin{proof}
\begin{enumerate}[label={$\langle1\rangle$\arabic*.},leftmargin=*]
\item The section identifies $C$ with the positive-genus curve $B$.
The proper nonconstant map $C\to D$ is finite. At a general point of
$D$, the cusp normalization
\[
 (u,z)\longmapsto(u^3,u^2,z)
\]
has one branch and degree one on its singular curve. Thus $C\to D$ is
the normalization of $D$.
Proposition~\ref{cdp:st:A.5.14}, applied to the same $Y,S,D$, makes
$D$ smooth elliptic with constant cusp transversal. It follows that
$C\to D$ is an isomorphism and $B$ is elliptic.
Let $C_T\subset T$ be the inverse-image curve.
Proposition~\ref{cdp:st:C2d} gives a smooth neighbourhood of $C_T$,
the isomorphism $C_T\to D$, and
$\deg N_{C_T/T}=0$. Its leading normal terms identify the square and
cube of this normal line with the radical and Hessian quotient lines.

\item Factor the given resolution as
\[
 Z\xrightarrow{\sigma}Z'\xrightarrow{\mu}T,
 \qquad \pi=\mu\sigma,
\]
where $\mu$ is the minimal resolution and $\sigma$ is a sequence of
point blowups. Every $\sigma$-exceptional curve is $\pi$-exceptional
and is vertical for $p$. The map $p$ therefore descends to
$p':Z'\to B$. Explicitly, the proper image of $(\sigma,p)$ in
$Z'\times B$ projects finitely and bimeromorphically to the normal
surface $Z'$ with one point in every fibre. This projection is an
isomorphism and defines $p'$, with $p=p'\sigma$.
The curve $C'=\sigma(C)$ is a section and $\sigma|_C$ is an isomorphism.
Since $\mu$ is an isomorphism near $C_T$, one has
$(C')^2=\deg N_{C_T/T}=0$.

The strict transform of every $\mu$-exceptional curve is
$\pi$-exceptional, so those curves are vertical for $p'$.
Proposition~\ref{cdp:st:A.5.3} supplies the canonical effective divisor
on this minimal resolution. Its nonexceptional components have the
same conductor valuations as on $Z$. Consequently it has the form
$2C'+V'$, with $V'$ effective and vertical, and
\[
 \omega_{Z'}(2C'+V')\simeq(\eta\mu)^*\omega_S,
 \qquad c_1(K_{Z'}+2C'+V')=0\quad\hbox{in }H^2(Z',\Q).
\]
The constant cusp charts and the isomorphism of $\mu$ there show that
$V'$ is disjoint from $C'$.

\item Contract the vertical $(-1)$-curves of the irrational ruled
surface to obtain $\rho:Z'\to Z_0$, a $\mathbf P^1$-bundle over $B$.
Every rational curve is vertical because $B$ has positive genus.
Write $C_0=\rho(C')$ and $V_0=\rho_*V'$.
The image $C_0$ remains a section. The point-blowup formula for its
self-intersection gives $C_0^2=(C')^2+n=n\ge0$, where $n$ counts the
centres lying on the successive strict transforms of the section.
The divisor $V_0$ is an effective sum of fibres.
Pushing the canonical equality forward and intersecting with $C_0$
gives, by adjunction on the elliptic section,
\[
 0=(K_{Z_0}+2C_0+V_0)\cdot C_0
   =C_0^2+V_0\cdot C_0=n+V_0\cdot C_0.
\]
Thus $n=0$ and $V_0=0$.

For completeness, the conductor transform at a point blowup
$\beta:X\to X_0$ is determined by
$K_X=\beta^*K_{X_0}+E$: its exceptional coefficient is
$\operatorname{mult}_p(D_0)-1$. This is the same canonical divisor
identity used above. A ruling $(-1)$-curve meeting $C'$ cannot be a
component of $V'$, since $C'\cap V'=\varnothing$; its intersection
with $K_{Z'}+2C'+V'$ would be at least $-1+2=1$.

Since $n=0$, one has $C'=\rho^*C_0$.
All components of $V'$ are $\rho$-exceptional, and
\[
 K_{Z'}=\rho^*K_{Z_0}+E_\rho
\]
with $E_\rho$ effective and positive on every exceptional prime.
Hence $E_\rho+V'$ is numerically trivial and exceptional.
The total transforms of the successive exceptional curves of a
point-blowup sequence give an orthogonal basis with square $-1$.
The exceptional intersection form is therefore negative definite.
It follows that $E_\rho+V'=0$. Effectiveness gives
$E_\rho=V'=0$, so $\rho$ is an isomorphism and $Z'$ is the minimal
ruled bundle.

\item A vertical irreducible curve of this bundle is a fibre and has
square zero. To compare it with the exceptional curves of $\mu$,
use the local-function form of the intersection argument in
\cite[\S1, p.~6]{Mumford1961}.
For an exceptional fibre over a normal point, choose a nonzero
holomorphic function $g$ vanishing at that point. On the resolution
its divisor is $G+\sum m_iE_i$, with $m_i>0$ and $G$ the strict
transform of the zero curve. Pairing this principal divisor with
$E_j$ gives
$\sum_i m_i(E_i\cdot E_j)=-G\cdot E_j\le0$, with a strict inequality
for some $j$. The strict transform meets the exceptional fibre by
properness, and that fibre is connected because the target is normal.
For $A_{ij}=m_im_j(E_i\cdot E_j)$, its off-diagonal entries are
nonnegative and
\[
 x^{\mathsf T}Ax
 =\sum_i\left(\sum_jA_{ij}\right)x_i^2
     -\sum_{i<j}A_{ij}(x_i-x_j)^2.
\]
A vector in its zero locus has zero coordinate at a strict row, and
connectedness forces all coordinates to vanish. The exceptional
intersection form is negative definite. This argument uses the
displayed local holomorphic divisor and applies to the analytic
surface $T$.
Thus a square-zero fibre cannot be $\mu$-exceptional. The map $\mu$
has no exceptional curves, hence is finite and bimeromorphic; the
normality of $T$ makes it an isomorphism. Therefore $T=Z'$ is smooth
and minimal ruled over the elliptic curve $B$.

\item Finite duality on the smooth normalization identifies the
conductor ideal with the invertible ideal of its divisor, now
$2C_T$. Thus $\eta$ is an isomorphism off $C_T$, and its restriction
$C_T\to D$ is an isomorphism. The normalization is consequently a
proper continuous bijection and a homeomorphism of the underlying
compact spaces. Equivalently, additivity of Euler characteristic
gives
\[
 e(S)=e(T)-e(C_T)+e(D)=0.
\]
Here $e(T)=e(\mathbf P^1)e(B)=2\cdot0$, and both curves have genus one.
The full scheme fibre $mS$ has the same underlying space as $S$,
so its topological Euler characteristic is also zero.
\end{enumerate}
\end{proof}